\section{Introduction}
\label{sec:intro}

In July 2025 the Bank of England published, for the first time, a full technical description of the Bayesian structural vector autoregression (SVAR) it uses in the monetary policy round to separate domestic from global drivers of the UK business cycle \citep{brignone2025}. The model is small --- eight quarterly variables, six identified structural shocks --- but it carries substantial weight in the Bank's process: it supplies impulse responses, forecast-error variance decompositions, real-time structural shock readings, and a round-by-round structural narrative for revisions to the Monetary Policy Report forecast. No official replication package accompanies the paper.

This paper documents an independent, open-source Python replication of that model, developed as part of \emph{PolicyEngine Macro}, PolicyEngine's suite of open macroeconomic models for the UK and US \citep{policyengine}.\footnote{Code: \url{https://github.com/PolicyEngine/boe-var-model}; hosted interface: \url{https://policyengine-macro.vercel.app/svar/}.} Within the suite the SVAR plays a specific role: it is the \emph{baseline and conditioning member}. It reads the current state of the economy in structural-shock terms, produces unconditional forecasts with credible bands, and generates revision narratives that other models in the suite condition on. It deliberately does \emph{not} score policy reforms --- reform scoring is the province of the suite's overlapping-generations, macroeconometric and microsimulation members --- because a sign-restricted SVAR identifies shocks, not policy rules.

The contribution is threefold. First, \emph{replication as a public good}: the Bank released no code or data, and its two world aggregates are internal, unpublished series. Reconstructing the pipeline --- data assembly, Minnesota-prior Bayesian estimation, zero-and-sign-restriction identification following \citet{arias2018}, and the standard outputs --- makes the Bank's analytical machinery inspectable and extensible by anyone. Second, \emph{validation with numbers rather than adjectives}: Section~\ref{sec:validation} compares the replication against the paper's published figures quantitatively, reporting exact satisfaction of every identifying restriction, the sign of key unrestricted responses, and variance-decomposition shares side by side with the paper's, together with machine-precision internal-consistency checks. Third, \emph{operationalisation}: the model runs as a set of hosted Model Context Protocol (MCP) tools --- \texttt{forecast\_uk}, \texttt{latest\_shocks}, \texttt{model\_summary} --- so that forecasts and structural readings are callable from AI agents and ordinary scripts alike.

Section~\ref{sec:literature} places the model in the Bayesian-VAR, sign-restriction and central-bank-suite literatures. Section~\ref{sec:model} sets out the model, the prior and posterior algebra, the identification algorithm with its importance weights, and the mathematics of the impulse-response, variance-decomposition, historical-decomposition and forecasting outputs. Section~\ref{sec:source} identifies the source paper and delimits the replication scope. Section~\ref{sec:implementation} describes the implementation. Section~\ref{sec:validation} reports the calibration and validation results together with the replication's core structural figures. Section~\ref{sec:results} shows a forecast from the 2024Q2 data edge and confronts it, quarter by quarter, with subsequent ONS outturns and the Bank's own contemporaneous and later Monetary Policy Report narratives. Section~\ref{sec:limitations} concludes with limitations. Appendices document the full dataset construction, the prior hyperparameters, and the sampling diagnostics.
